\chapter{MARCO METODOLÓGICO}

Este capítulo describe el método completo, organizado en las cuatro
unidades de la arquitectura. Comienza por la construcción del entorno
de simulación (U1), con dos caminos posibles: la importación de una
intersección real desde OpenStreetMap, prevista para el cruce de Lima,
y la descripción directa de nodos y vías, que es la que se usó para los
escenarios evaluados en esta fase. Sigue con la formulación e
implementación del agente de aprendizaje por refuerzo (U2), la
instrumentación de las métricas (U3) y el plan de pruebas con el que se
obtienen los resultados del Capítulo~\ref{cap:resultados}. La
coordinación distribuida entre semáforos vecinos (U4), formulada en el
Capítulo~I, se implementa en la Sección~\ref{sec:u4} sobre los
escenarios de varios semáforos.


\section{Alcance y unidades}

\noindent\textbf{U1: Preprocesamiento.} Comprende todas las tareas
relacionadas con la construcción del escenario: la instalación y
configuración de SUMO, la obtención de la red (por conversión de un
recorte de OSM o por descripción directa de nodos y vías), la
definición de la demanda vehicular y la creación de archivos de
configuración reproducibles (\texttt{.net.xml} para la red,
\texttt{.rou.xml} para la demanda y \texttt{.sumocfg} para el
escenario). El resultado de U1 es un escenario listo para simular.

\noindent\textbf{U2: Entrenamiento.} Toma como entrada el escenario de
U1 y ejecuta SUMO a través de la interfaz TraCI para que un agente de
aprendizaje por refuerzo observe el estado del tráfico y modifique
las fases de los semáforos.  En U2 no se modifican la red ni las
rutas: el enfoque se centra en el algoritmo de control.

\noindent\textbf{U3: Medición.} Calcula, sobre las mismas salidas de
SUMO y para todos los controles comparados, las métricas de desempeño
con las que se evalúa el sistema (Sección~\ref{sec:u3}).

\noindent\textbf{U4: Coordinación distribuida.} Añade mensajería entre
agentes vecinos, estado extendido, penalización por \textit{spillback}
y un algoritmo con aproximación de funciones (Sección~\ref{sec:u4}).
Se evalúa en variantes incrementales sobre los escenarios de varios
semáforos, frente a los mismos controles de referencia que U2, de modo
que el aporte de la comunicación quede aislado del resto de cambios.


\section{Resumen de las unidades}

\begin{description}
  \item[U1:] Instalar y configurar SUMO; importar la intersección
    mediante el asistente OSM o descarga manual; limpiar y validar
    la red; generar la demanda vehicular y el archivo de configuración
    \texttt{.sumocfg}.  El resultado es un escenario reproducible.
  \item[U2:] Ejecutar el escenario con TraCI y entrenar un agente de
    aprendizaje por refuerzo que ajuste las fases semafóricas.  U2
    depende exclusivamente de los archivos generados en U1 y no
    modifica la red ni las rutas.
  \item[U3:] Medir el desempeño de cada control con las mismas salidas
    de SUMO (viajes por vehículo, colas por TraCI, emisiones) y
    compararlo con el protocolo del plan de pruebas
    (Secciones~\ref{sec:u3} y~\ref{sec:plan}).
  \item[U4:] Coordinar los agentes de semáforos vecinos: estado
    continuo con la información del vecino, mensajería entre cruces,
    recompensa con penalización por \textit{spillback} y aprendizaje
    con IPPO (Sección~\ref{sec:u4}).
\end{description}

%----------------------------------------------------------------------

\section{Escenarios sintéticos de prueba}
\label{sec:escenarios}

Mientras se completa la importación del cruce real de Lima, los
módulos U2, U3 y U4 se desarrollaron y validaron sobre siete escenarios
sintéticos, construidos con \texttt{netconvert} a partir de archivos de
nodos y vías escritos a mano, del más simple al más complejo
(Tabla~\ref{tab:escenarios}). Los cuatro primeros comparten el mismo
nivel de demanda; el quinto, corredor\_alta, es el escenario de estrés
del módulo U4 y se describe en la Sección~\ref{sec:u4}; los dos
últimos, red\_alta y malla3, son los casos P4 y P5 del plan de pruebas
y se describen en la Sección~\ref{sec:plan2}. Todos comparten el mismo tipo de
vehículo (turismo de 5~m, aceleración 2.6~m/s$^2$, velocidad máxima
50~km/h, imperfección del conductor $\sigma = 0.5$) y el mismo patrón
de demanda: una hora dividida en dos periodos de 30 minutos con la
punta invertida, de modo que un programa de tiempo fijo no pueda ser
óptimo en ambos a la vez. Los flujos se definen con
\verb|<flow ... vehsPerHour>| por movimiento y periodo, lo que produce
llegadas equiespaciadas; el ámbar es de 3~s en todos los casos (valor
por defecto de \texttt{netconvert}).

\begin{table}[H]
\setstretch{1}
\centering
\footnotesize
\setlength{\tabcolsep}{3pt}
\begin{tabular}{p{1.5cm}p{2.7cm}p{1.8cm}p{2.9cm}p{1.0cm}p{2.0cm}}
\toprule
Escenario & Geometría & Fases verdes & Demanda del periodo punta (veh/h) & Vehículos en la hora & Programa fijo de referencia \\
\midrule
cruce & cruce N--S / E--O, 1 carril por sentido & 2 & eje principal 650 por sentido; transversal 150; giros 40 a 80 & 1844 & verde 25~s, ciclo 56~s \\
cruce2 & 2 carriles por sentido más carril exclusivo de giro a la izquierda & 4 (2 principales y 2 de giro protegido) & recto 420 por sentido; giro izquierdo 160; transversal 200 & 1884 & verde 20~s, giro 6~s, ciclo 64~s \\
corredor & dos cruces separados 300~m sobre una avenida E--O, 1 carril por sentido & 2 por cruce & avenida 600 en el sentido punta y 250 en el contrario; transversales 180 & 1810 & avenida 20~s, transversales 10~s, ciclo 36~s; desfase de 15~s en el segundo cruce \\
red & rotonda sin semáforo más tres cruces separados 300~m sobre una avenida de 2 carriles por sentido & 2 por cruce & avenida 650 y 250; transversales 150; rotonda 80 a 150 & 2152 & avenida 25~s, transversales 15~s, ciclo 46~s; desfases 0, 23 y 0~s \\
corredor\_alta & la red de corredor con toda la demanda multiplicada por 2 (escenario de estrés del módulo U4) & 2 por cruce & avenida 1200 y 500; transversales 360; giros 120 & 3620 & verde 25~s en la avenida y 15~s en las transversales; desfase de 25~s \\
red\_alta & la red de red con toda la demanda multiplicada por 1.2 (caso P4) & 2 por cruce & avenida 780 y 300; transversales 180; rotonda 96 a 180 & 2582 & avenida 30~s, transversales 10~s, ciclo 46~s; desfases 0, 20 y 0~s \\
malla3 & malla de 3 $\times$ 3 cruces separados 300~m, 2 carriles por sentido en todas las vías (caso P5) & 2 por cruce & por fila 600 y 250; por columna 180 por sentido; 12 giros de 60 & 4350 & avenida 25~s, transversales 15~s, ciclo 46~s; desfases en ajedrez, 0 y 23~s \\
\bottomrule
\end{tabular}
\caption{Escenarios sintéticos de prueba. La demanda del segundo periodo es la del primero con la punta invertida.}
\label{tab:escenarios}
\end{table}
% fuente (red_alta, malla3): plan2/red_alta.rou.xml, malla3.rou.xml; programa fijo: plan2/<esc>.add.xml;
% vehiculos en la hora de corredor_alta, red_alta y malla3: columna vehiculos del brazo fijo en plan2/evaluacion_<esc>.csv


Como segunda referencia se usa el control actuado nativo de SUMO, con
las mismas fases y los mismos límites de verde mínimo y máximo que el
agente (Sección~\ref{sec:plan}).


\section{Desarrollo del módulo U2: Aprendizaje por Refuerzo para el control de semáforos}

El segundo módulo de la metodología aborda el diseño e
implementación de un controlador inteligente que modifique las fases
de los semáforos en función del estado del tráfico. Para ello se
recurre a técnicas de \emph{Aprendizaje por Refuerzo} (\emph{Reinforcement
Learning}, RL), en las que un agente aprende a tomar decisiones
mediante la interacción con un entorno simulado. Este apartado
describe la formulación del problema, la interfaz con SUMO y los
algoritmos de RL más habituales para el control de intersecciones.

\subsection*{Formulación del problema}

En RL se define un \emph{entorno} compuesto por un conjunto de
\emph{estados}, un conjunto de \emph{acciones} y una función de
\emph{recompensa}. La formulación implementada en esta tesis para un
cruce con cuatro aproximaciones es la siguiente:

\begin{itemize}
  \item \textbf{Estado:} el número de vehículos detenidos en cada
    aproximación, leído con \verb|edge.getLastStepHaltingNumber| y
    discretizado en cuatro niveles (0; 1 a 3; 4 a 8; más de 8), más el
    índice de la fase verde activa. El estado es $s = (c_N, c_S, c_E,
    c_O, f)$ con $c_i \in \{0,1,2,3\}$. Con dos fases verdes el espacio
    tiene $4^4 \times 2 = 512$ estados y con cuatro fases, $1024$. Esta
    discretización es la que hace viable una tabla Q y, a la vez, la
    que limita al método tabular cuando crece el número de fases, como
    se muestra en el Capítulo~\ref{cap:resultados}.

  \item \textbf{Acción:} dos acciones. $a = 0$ mantiene la fase verde
    actual y $a = 1$ pasa a la siguiente fase verde del programa,
    insertando automáticamente el ámbar del programa (3~s). El agente
    no elige una fase arbitraria ni una duración: decide únicamente
    cuándo termina el verde en curso.

  \item \textbf{Recompensa:} negativa y proporcional a la congestión en
    el instante de la decisión,
    \begin{equation}
      r = -\Bigl(w_1 \sum_i c_i + w_2 \sum_i e_i\Bigr), \qquad
      w_1 = 1.0,\; w_2 = 0.01,
      \label{eq:recompensa}
    \end{equation}
    donde $c_i$ es el número de vehículos detenidos en la aproximación
    $i$ (sin discretizar) y $e_i$ es la suma de los tiempos de espera de
    los vehículos que en ese momento están en la aproximación $i$
    (\verb|edge.getWaitingTime|). Con estos pesos domina el término de
    colas; el de espera desempata entre estados con la misma cola. La
    recompensa mide el nivel de congestión, no su incremento entre
    decisiones. Es la recompensa de la fase exploratoria. En el plan de
    pruebas vigente (Sección~\ref{sec:plan2}) Q-learning y SARSA se
    entrenan con la recompensa de retraso de la
    Ecuación~\ref{eq:recompensa-u4}, la misma de DQN, IPPO y MAPPO; la
    de la Ecuación~\ref{eq:recompensa} se conserva solo en la ablación
    \texttt{ql\_colas}.

  \item \textbf{Restricciones operativas:} el verde dura al menos 10~s
    antes de que el agente pueda cambiarlo (5~s en las fases de giro
    protegido, identificadas por tener a lo más dos movimientos en
    verde) y como máximo 60~s, tras los cuales el cambio se fuerza. El
    agente decide cada 5~s.
\end{itemize}

Con estos elementos se modela el control del semáforo como un
\emph{problema de decisión de Markov} (MDP). El agente observa el
estado, aplica una acción, avanza la simulación y recibe una
recompensa; a partir de la secuencia de experiencias va aprendiendo
una política de control. En los escenarios con varios cruces se
instancia un agente por semáforo, cada uno con su propia tabla Q, su
estado local y su recompensa local; los agentes no intercambian
información (eso es lo que añade el módulo U4).


\subsection*{Algoritmos recomendados}

Diversas técnicas de RL han sido aplicadas al control de
semáforos. Los algoritmos tabulares, como \emph{Q-learning} o
\emph{SARSA}, son adecuados para intersecciones sencillas con un
número limitado de estados, y son los que permiten inspeccionar
directamente la política aprendida. Para espacios de estado mayores se
emplean métodos de \emph{Deep Reinforcement Learning} como \emph{Deep
Q-Network} (DQN) y \emph{Proximal Policy Optimization} (PPO), que
sustituyen la tabla por una red neuronal capaz de generalizar entre
estados parecidos. Ambos tienen implementación de referencia en
\texttt{stable-baselines3}, sobre PyTorch, y su acoplamiento a SUMO
reutiliza el mismo bucle de TraCI descrito arriba.

En esta fase se implementaron los cuatro, sin
\texttt{stable-baselines3}: Q-learning y SARSA tabulares (este módulo,
\texttt{rl\_corredor.py}), PPO en su versión multiagente independiente
(IPPO, módulo U4, Sección~\ref{sec:u4}) y, para el plan de pruebas, un
Double DQN y un MAPPO en PyTorch puro (Sección~\ref{sec:plan2}).
La motivación para pasar de la tabla a la red está en los resultados
del Capítulo~\ref{cap:resultados}: la representación tabular alcanza su
techo en cuanto crece el número de fases del semáforo, que es
exactamente el caso en que la aproximación por redes neuronales resulta
pertinente.



%----------------------------------------------------------------------
\section{Módulo U3: medición del desempeño}
\label{sec:u3}

Todas las métricas se calculan con las mismas salidas de SUMO para los
tres controles comparados (tiempo fijo, actuado y agente RL), de modo
que la comparación no dependa de cómo decide cada uno.

\begin{itemize}
  \item \textbf{Por vehículo}, del archivo \texttt{tripinfo} que SUMO
    escribe al terminar cada viaje: el retraso (\texttt{timeLoss}), la
    espera detenida (\texttt{waitingTime}), el número de paradas
    (\texttt{waitingCount}, veces que la velocidad bajó de 0.1~m/s) y
    el instante de llegada. El \textit{throughput} se define como el
    número de vehículos que completan su viaje antes de los 3600~s. Con
    el dispositivo de emisiones activado
    (\verb|--device.emissions.probability 1|) cada viaje incluye además
    sus emisiones acumuladas; se reportan CO$_2$ en gramos y NO$_x$ en
    miligramos por vehículo, con la clase de emisión por defecto de SUMO
    (HBEFA3, turismo a gasolina Euro~4) \citep{SUMO_Emissions}. Activar
    el dispositivo no altera la dinámica ni la secuencia aleatoria de la
    simulación: las métricas con y sin él coinciden.
  \item \textbf{Por segundo}, vía TraCI: el número de vehículos
    detenidos, sumado sobre todas las aproximaciones semaforizadas del
    escenario. De esa serie salen la cola promedio, que es su media por
    segundo de simulación, y la cola máxima, que es su mayor valor
    instantáneo; se guarda además un muestreo por minuto para observar
    la respuesta de cada control al cambio de demanda.
  \item \textbf{Por decisión}: la recompensa de la
    Ecuación~\ref{eq:recompensa} se evalúa cada 5~s de reloj fijo, con
    el mismo muestreo bajo los tres controles, para comparar lo que el
    agente optimiza con lo que consiguen el fijo y el actuado. Se
    reporta como promedio por decisión y no como suma acumulada, para
    que el valor no dependa de la duración del episodio. Se cuentan
    además los cambios de fase (entradas a una fase verde) de cada
    control, sobre el episodio completo.
  \item \textbf{Retornos por segundo}, desde el plan de pruebas
    vigente: tres recompensas (de retraso, de colas y de
    \textit{spillback}) acumuladas cada segundo sobre todos los
    semáforos, con cualquier control. Se definen en la
    Sección~\ref{sec:plan2}.
\end{itemize}

\subsection*{Métricas de corredor y \textit{spillback}}

En los escenarios de varios cruces se miden además, con las mismas
salidas y para todos los controles, las magnitudes que la coordinación
debería mover. Del \texttt{tripinfo}, sobre los vehículos que recorren
la avenida de extremo a extremo (850 en corredor, unos 900 en red): el
tiempo de viaje y el retraso medios, el retraso medio en el sentido de
la punta (el que favorece la onda verde del programa fijo), las paradas
medias y la fracción de vehículos que se detiene a lo sumo una vez,
que es el indicador más directo de progresión. Por TraCI, en cada
carril de los enlaces que unen dos cruces: el máximo de vehículos
detenidos y el máximo de vehículos presentes durante la corrida, y el
\textit{spillback}, definido como un segundo en el que algún carril de
algún enlace tiene detenidos al menos el 90\,\% de sus plazas (34 de
38: la cola queda a menos de 30~m del cruce aguas arriba y el siguiente
vehículo que se detenga bloquea la intersección). Se reportan la
fracción de segundos de la hora con \textit{spillback}, los segundos
totales y el número de episodios de \textit{spillback} (flancos de
subida). Se cuenta también el número de vehículos que no terminaron su
viaje al cerrar la simulación, que es cero con la demanda nominal. Al
pedir la serie se guarda además la fase de cada semáforo en cada
segundo, para observar el desfase que emerge entre cruces.

Dos precisiones importan para leer los resultados. El retraso de SUMO
incluye tiempo perdido que no depende del semáforo (arranque desde
velocidad cero, giros a menor velocidad, vacilaciones del conductor),
así que un cambio porcentual sobre \texttt{timeLoss} subestima el
efecto sobre el retraso atribuible al control. Y la espera detenida es
la señal que optimiza la recompensa, pero no es el retraso: un control
que fragmenta las colas en paradas más cortas puede bajar la espera y
subir las paradas sin mover el retraso. Por eso el retraso es la
métrica primaria y la espera, las colas, las paradas, el
\textit{throughput} y las emisiones son secundarias.

%----------------------------------------------------------------------

\section{Desarrollo del módulo U4: coordinación distribuida entre semáforos}
\label{sec:u4}

Los resultados del módulo U2 en los escenarios de varios cruces
(Capítulo~\ref{cap:resultados}) motivan este módulo: tres agentes
tabulares independientes pierden un $15$\,\% de retraso frente a la onda
verde tuneada de red, dos agentes en el corredor pierden un $60$\,\%
frente al fijo con reparto de verde tuneado (un $6$\,\% frente al fijo
simétrico de la primera ronda) y su política resultó sensible a un
cambio de 2~s en el verde mínimo. El
módulo U4 cambia cuatro cosas a la vez, de forma incremental y
comparable: el estado pasa de bins a un vector continuo que incluye lo
que el propio detector ve en movimiento; cada semáforo recibe mensajes
de sus vecinos sobre la avenida; la recompensa puede penalizar la cola
que el agente empuja sobre el enlace que comparte con el vecino
(\textit{spillback}); y el aprendizaje pasa de la tabla Q a
\emph{Independent PPO} (IPPO) con una sola red compartida por todos los
semáforos. Todo se implementa en \texttt{rl\_ippo.py} (Anexo~B) sobre
PyTorch en CPU, sin Gym ni \texttt{stable-baselines3}, reutilizando el
bucle de TraCI, el reloj de decisión y las métricas de U2 y U3.

\subsection*{Vecinos y canal de mensajes}

Dos semáforos son vecinos cuando comparten un tramo de la avenida: en
corredor C1 y C2; en red C1--C2 y C2--C3 (la rotonda no tiene semáforo
y no participa). En la malla de nueve cruces del plan de pruebas (P5)
cada semáforo tiene de dos a cuatro vecinos, uno por lado
(Sección~\ref{sec:plan2}). Cada 5~s de reloj global, cada semáforo publica a cada
vecino un mensaje con cinco datos que un controlador real podría
transmitir: la fase en curso (avenida, transversal o ámbar), el tiempo
que lleva en ella, los vehículos presentes en la aproximación que
descarga hacia el vecino, la cola detenida en esa aproximación y la
cola detenida en el enlace que une al emisor con el receptor, que es
el enlace por el que llegan al receptor esos vehículos. El receptor
conserva el último mensaje recibido y su edad. El canal admite una
latencia fija y una probabilidad de pérdida por mensaje, ambas nulas
durante el entrenamiento y usadas solo en la prueba de robustez del
plan de pruebas; la pérdida usa un generador aleatorio propio para no
alterar la secuencia de la simulación.

\subsection*{Observación, acciones y máscara}

La observación de cada agente es un vector de 31 componentes
normalizadas entre 0 y 1 (Tabla~\ref{tab:obs-u4}): un bloque local de
15 y dos ranuras fijas de 8, una para el vecino del oeste y otra para
el del este. Las colas se dividen por la capacidad de la aproximación
(plazas de 7.5~m por carril) y la cola del enlace por sus plazas
(38 en los enlaces de 300~m). La variante \emph{local} deja las dos
ranuras a cero: contiene todo lo que un detector propio puede medir, de
modo que la diferencia entre \emph{vecinos} y \emph{local} se atribuye
solo a la mensajería. La misma red y la misma disposición sirven para
los tres escenarios; el identificador del semáforo ocupa tres
posiciones fijas. En la malla del plan de pruebas (P5) la disposición
se extiende sin cambiar el bloque local: nueve posiciones de
identificador y cuatro ranuras de 8 (oeste, este, norte y sur), 53
componentes en total; ``verde hacia mí'' es el verde este--oeste para
un vecino al oeste o al este y el norte--sur para uno al norte o al
sur. Con dos lados la disposición sigue siendo la de la tabla, y las
políticas anteriores cargan sin cambios.

\begin{table}[H]
\centering
\scriptsize
\setlength{\tabcolsep}{3pt}
\begin{spacing}{1}
\begin{tabular}{llp{9.2cm}}
\toprule
Bloque & Componentes & Contenido \\
\midrule
Local & 1--4 & cola detenida en cada aproximación (O, E, N, S) sobre su capacidad \\
 & 5--8 & vehículos en movimiento en cada aproximación sobre su capacidad \\
 & 9--10 & fase verde activa (uno de dos) \\
 & 11 & tiempo en el verde actual sobre 60~s \\
 & 12 & 1 si ya se cumplió el verde mínimo \\
 & 13--15 & identificador del semáforo (uno de tres) \\
\midrule
Vecino oeste & 16 & 1 si existe el vecino \\
(17--22 a cero si & 17 & 1 si el vecino está en la fase de la avenida (descarga hacia mí) \\
no hay mensaje) & 18 & 1 si el vecino está en ámbar \\
 & 19 & tiempo del vecino en su fase sobre 60~s \\
 & 20 & vehículos en la aproximación del vecino que descarga hacia mí, sobre su capacidad; cero si el vecino no está en la fase de la avenida \\
 & 21 & cola detenida en esa aproximación sobre su capacidad \\
 & 22 & cola detenida en el enlace que llega desde el vecino (del vecino hacia mí), máximo por carril, sobre sus plazas \\
 & 23 & edad del último mensaje sobre 30~s (1 si nunca llegó) \\
\midrule
Vecino este & 24--31 & igual que 16--23 \\
\bottomrule
\end{tabular}
\end{spacing}
\caption{Observación de cada agente del módulo U4.}
\label{tab:obs-u4}
\end{table}



Las acciones son las de U2: mantener el verde o pasar a la siguiente
fase con ámbar automático. El verde mínimo y máximo se imponen con una
máscara sobre las acciones: antes del mínimo solo es válida mantener,
a los 60~s solo cambiar, y entre ambos las dos. Las decisiones forzadas
por la máscara no aportan gradiente a la política (su probabilidad es
1) pero sí al crítico. Se reportan por episodio las decisiones totales,
comparables con U2, y las decisiones libres.

\subsection*{Recompensa y variantes}

La recompensa de U4 conserva la estructura y los pesos de la
Ecuación~\ref{eq:recompensa} ($w_1 = 1$, $w_2 = 0.01$) con dos cambios
que salieron de la prueba de humo del plan de pruebas (más abajo). El
término de cola detenida se sustituye por los vehículos-equivalentes
retrasados de cada aproximación, $n_a\,(1 - \bar v_a / v_{\max,a})$,
donde $n_a$ es el número de vehículos presentes y $\bar v_a$ su
velocidad media: un vehículo detenido cuenta 1 y uno que avanza a
media velocidad cuenta 0.5. Es la versión por segundo del
\texttt{timeLoss} de SUMO, la métrica primaria, y se lee de la misma
suscripción por aproximación. Y el agente recibe, por cada transición,
la suma por segundo de esa cantidad sobre su intervalo de decisión
dividida entre 5~s (retraso-segundos), en lugar del nivel en el
siguiente instante de decisión: en un intervalo de mantener coincide
con el nivel medio y en uno de cambiar, de 8~s, pesa $8/5$. La variante
\emph{spill} le resta además un término por la cola que el agente
empuja sobre cada enlace que alimenta hacia un vecino, acumulado de la
misma forma:
\begin{equation}
  \begin{aligned}
  r_i = -\frac{1}{5}\sum_{t \in \text{intervalo}} \Bigl[\;
        & w_1 \sum_a n_a(t)\Bigl(1 - \frac{\bar v_a(t)}{v_{\max,a}}\Bigr)
        + w_2 \sum_a \text{espera}_a(t) \\
        & + w_3 \sum_{e \in \text{salidas}_i} \max\bigl(0,\; q_e(t) - P_e/2\bigr) \Bigr],
  \end{aligned}
  \label{eq:recompensa-u4}
\end{equation}
donde $q_e$ es el máximo por carril de vehículos detenidos en el enlace
$e$, $P_e$ sus plazas por carril (38; 37 en malla3) y $w_3 = 2$ (cero en \emph{local} y
\emph{vecinos}): el agente aguas arriba ve la cola que bloquea al
vecino como si fuera propia, con peso doble, desde que ocupa la mitad
del enlace. La variante \emph{cambio}, opcional, resta $w_4 = 2$ por
cada cambio de fase ejecutado, elegido o forzado, y se corre aparte
porque cambia el objetivo. Los pesos están fijados de antemano, sin
barrido. Lo que ve el algoritmo es la recompensa multiplicada por
0.02; en el registro por episodio se conserva la recompensa de U2
(colas detenidas) a reloj fijo de 5~s, comparable con los cuatro casos
ya evaluados, y aparte la recompensa completa que vio el agente.



%----------------------------------------------------------------------
\section{Plan de pruebas}
\label{sec:plan}
\label{sec:plan2}

El plan de pruebas vigente, aprobado por el asesor el 23.09.2026,
aplica un mismo método en todas las configuraciones y mide las mismas
salidas en todas ellas. Cada caso fija un escenario (red y demanda);
en cada caso se entrenan y evalúan los mismos modelos con el mismo
protocolo, y un veredicto prefijado dice si el modelo funciona en ese
caso. La fase anterior, exploratoria, se resume al final de la
sección.


\subsection*{Fase exploratoria previa}

Antes de este plan hubo una fase exploratoria, cuyos resultados
reporta el Capítulo~\ref{cap:resultados}, con otra numeración de los
casos: E1 cruce, E2 cruce2, E3 corredor y E4 red para el agente
tabular, E3-U4 y E4-U4 para las variantes IPPO en corredor y red, y E6
para corredor\_alta. En ella el Q-learning se entrenaba con una sola
semilla base, entre 40 y 120 episodios y con la recompensa de colas;
IPPO, con tres bases de 200 episodios y la recompensa de retraso; en
los dos casos se publicaba la política del último episodio, y la
evaluación usaba diez
semillas (1001 a 1010), con una validación intermedia en la semilla
999. El éxito de U4 se juzgaba con siete niveles prefijados, de (N0)
medición (regresión, determinismo y prueba de humo) a (N6) la meta de
$-25$\,\% frente al fijo, pasando por estabilidad entre semillas base
(N1), ventaja sobre Q-learning (N2), coordinación (N3),
\textit{spillback} (N4) y paridad con el actuado (N5). Esa fase mostró
que el Q-learning con recompensa de colas quedaba por detrás del
actuado y que la última política no siempre era la mejor, pero sus
casos no eran comparables entre sí: cambiaban el número de episodios,
la recompensa y la regla de selección según el modelo, y los fijos
nominales se habían tuneado con semillas de evaluación. El plan de
esta sección corrige esas diferencias y reemplaza sus conclusiones.


\subsection*{Objetivo de cada prueba}

La pregunta es la misma en todas las celdas de la matriz caso por
modelo: \emph{¿el modelo reduce el retraso por vehículo frente al
tiempo fijo tuneado en ese caso, y no queda por detrás del actuado?}
Cada caso agrega una dificultad a la anterior. P1 y P2 comparan un
corredor de dos semáforos con demanda nominal y con demanda doble, y
responden si el modelo sigue funcionando cuando la demanda supera la
capacidad de un carril. P3 y P4 repiten la comparación en la red de
tres semáforos, dos carriles por sentido y una rotonda sin semáforo,
con la demanda nominal y con una un 20\,\% mayor (el límite que impone
la rotonda), y responden si funciona con más carriles y con un
elemento que no controla. P5 pasa a una malla de nueve semáforos en la
que los cruces tienen vecinos en los dos ejes, hasta cuatro, y
responde si el método escala a una red.

\subsection*{Casos}

La Tabla~\ref{tab:plan2-casos} resume los casos. P1 a P5 se ejecutan
en esta fase. P6 (la misma malla con demanda alta) y P7 (una red de
cuatro por cuatro cruces, con tres carriles por sentido y fase de giro
protegido) quedan definidos como trabajo siguiente: repiten el mismo
protocolo sobre escenarios de mayor tamaño y no cambian el método.
La avenida es siempre el eje este--oeste y la transversal, el
norte--sur; en la malla los dos ejes tienen dos carriles y la
convención se mantiene.

\begin{table}[H]
\setstretch{1}
\centering
\footnotesize
\setlength{\tabcolsep}{4pt}
\begin{tabular}{llcclc}
\toprule
Caso & Escenario & Semáforos & \makecell{Carriles por\\sentido (avenida)} & Demanda & \makecell{Vehículos\\en la hora} \\
\midrule
P1 & corredor       & 2              & 1 & nominal                & 1810 \\
P2 & corredor\_alta & 2              & 1 & 2.0 $\times$ nominal   & 3620 \\
P3 & red            & 3 + rotonda    & 2 & nominal                & 2150 \\
P4 & red\_alta      & 3 + rotonda    & 2 & 1.2 $\times$ nominal   & 2580 \\
P5 & malla3         & 9 (3$\times$3) & 2 & nominal                & 4350 \\
\midrule
P6 & malla 3$\times$3 & 9            & 2 & alta                   & por definir \\
P7 & red 4$\times$4 & 16             & \makecell{3, con giro\\protegido} & por definir & por definir \\
\bottomrule
\end{tabular}
\caption[Casos del plan de pruebas]{Casos del plan de pruebas. P1 a P5 se ejecutan en esta
fase; P6 y P7 quedan como trabajo siguiente. Carriles por sentido en la avenida. La demanda
nominal es la de cada escenario en la Tabla~\ref{tab:escenarios} y, en malla3, la descrita en
esta sección; el factor de P4 se calibró antes de entrenar (Tabla~\ref{tab:red-alta-factor}).
Vehículos en la hora: suma de los flujos declarados; por el redondeo de las llegadas
equiespaciadas SUMO genera 2152 en red y 2582 en red\_alta (Tabla~\ref{tab:escenarios}).}
\label{tab:plan2-casos}
\end{table}
% fuente: plan2/veredicto.csv (columna demanda_veh_h, suma de vehsPerHour de cada .rou.xml);
% vehiculos generados: columna vehiculos de plan2/evaluacion_<esc>.csv



\subsection*{Modelos}

Cada celda compara los modelos de la Tabla~\ref{tab:plan2-modelos}.
Los dos controles de referencia no aprenden: el tiempo fijo, con el
programa tuneado para esa celda, y el actuado nativo de SUMO
\citep{SUMO_TLS_Actuated}, con las mismas fases y verde de 10 a 60~s
(Sección~\ref{sec:escenarios}). Los aprendices son dos tabulares
(Q-learning y SARSA), cuatro de aproximación de funciones que aprenden
cada uno con la información de su propio semáforo (las tres variantes
IPPO de la Sección~\ref{sec:u4} y un Double DQN con pesos compartidos)
y uno con crítico centralizado (MAPPO). DQN y MAPPO usan la misma observación que \emph{vecinos}, de
modo que la comparación entre algoritmos no se mezcla con la
comparación entre observaciones. La fila \texttt{ql\_colas} es una
ablación: el mismo Q-learning entrenado con la recompensa original de
U2 (Ecuación~\ref{eq:recompensa}), para medir cuánto del resultado se
debe a la recompensa y cuánto al algoritmo. Entre \texttt{ql} y
\texttt{ql\_colas} solo cambia la recompensa: el algoritmo, el estado,
las acciones y el protocolo son los mismos, así que la diferencia entre
sus dos filas aísla el efecto de la recompensa. Los hiperparámetros de
DQN y MAPPO se fijaron antes de entrenar en los casos y se comprobaron
con una prueba de humo en el cruce de un semáforo, que no es un caso
del plan; quedan registrados en el archivo de metadatos de cada
política y en la Tabla~\ref{tab:plan2-hiper}.

\begin{table}[H]
\setstretch{1}
\centering
\footnotesize
\begin{tabular}{l>{\raggedright\arraybackslash}p{6.0cm}ll}
\toprule
Modelo & Descripción & \makecell{Recompensa de\\entrenamiento} & Casos \\
\midrule
fijo & tiempo fijo tuneado por barrido en cada celda & -- & P1 a P5 \\
actuado & actuado nativo de SUMO, verde mínimo 10~s y máximo 60~s & -- & P1 a P5 \\
ql & Q-learning tabular, un agente por semáforo & retraso & P1 a P5 \\
sarsa & SARSA tabular, un agente por semáforo & retraso & P1 a P5 \\
ql\_colas & Q-learning con la recompensa de U2 (ablación) & colas & P1 a P4 \\
ippo\_local & IPPO con pesos compartidos, sin mensajes & retraso & P1 a P5 \\
ippo\_vecinos & IPPO con mensajes entre vecinos & retraso & P1 a P5 \\
ippo\_spill & IPPO con mensajes y término de \textit{spillback} & \makecell[l]{retraso +\\\textit{spillback}} & \makecell[l]{P1 y P2; en P3\\a P5 = vecinos} \\
dqn & Double DQN con pesos compartidos, observación de ippo\_vecinos & retraso & P1 a P5 \\
mappo & actor de ippo\_vecinos con crítico centralizado & retraso & P1 a P5 \\
\bottomrule
\end{tabular}
\caption[Modelos del plan de pruebas]{Modelos que se comparan en cada caso. La recompensa de
retraso es la Ecuación~\ref{eq:recompensa-u4} con $w_3 = 0$; la de colas, la
Ecuación~\ref{eq:recompensa}; la de \textit{spillback}, la Ecuación~\ref{eq:recompensa-u4} con
$w_3 = 2$. En P3, P4 y P5 el término de \textit{spillback} nunca se activa e ippo\_spill
coincide con ippo\_vecinos; la tabla de veredicto lo reporta como ``= vecinos''. La ablación
ql\_colas no se programó en P5.}
\label{tab:plan2-modelos}
\end{table}



\subsection*{Protocolo común}

Todos los aprendices (Q-learning, SARSA, DQN, IPPO y MAPPO) se
entrenan con la misma recompensa, la de retraso de la
Ecuación~\ref{eq:recompensa-u4} sin el término de \textit{spillback}
($w_1 = 1$, $w_2 = 0.01$, $w_3 = 0$): cada segundo se acumula por
semáforo el retraso de sus aproximaciones más la espera ponderada, y
la recompensa de una decisión es esa suma desde la decisión anterior
dividida entre 5. Es la recompensa con la que se entrenó IPPO en la
fase anterior. Las redes neuronales (DQN, IPPO y MAPPO) la reciben
multiplicada por 0.02 y los tabulares, sin escalar. Para el Q-learning
es un cambio respecto de la fase anterior, en la que se entrenaba con
la recompensa de colas de la Ecuación~\ref{eq:recompensa}; esa
recompensa se conserva solo en la ablación \texttt{ql\_colas}. La otra
excepción es la variante \emph{spill}, que por definición añade el
término de \textit{spillback}.

Cada aprendiz se entrena 200 episodios con tres semillas base (42,
2042 y 4042); el episodio $k$ usa la semilla de SUMO base$+k$. Cada
diez episodios, y en el último, se valida la política congelada
(acción de mayor valor o de mayor probabilidad, sin exploración) con
las semillas 999 y 1999, y el valor de validación es la media del
retraso en las dos. La validación guarda y restaura el estado de los
generadores aleatorios de Python, NumPy y PyTorch, así que el
entrenamiento reproduce el mismo registro por episodio, dígito a
dígito, con o sin ella. En cada validación se guarda un
\textit{checkpoint}; la política publicada de un entrenamiento es el
\textit{checkpoint} con menor valor de validación (en empate, el más
tardío), y no la del último episodio. De las tres semillas base, la
que representa al modelo en la celda es la mediana por ese mismo
valor. Ninguna de las dos elecciones usa las semillas de evaluación.

Los tabulares usan $\alpha = 0.1$, $\gamma = 0.9$ y
$\varepsilon$ inicial 1.0 con decaimiento geométrico por episodio
$0.01^{1/(0.8 N)}$, que para $N = 200$ episodios es 0.9716 y lleva
$\varepsilon$ a 0.01 al 80\,\% del entrenamiento, con 0.01 como piso.
IPPO conserva los hiperparámetros de la Tabla~\ref{tab:hiper-ippo};
los de DQN y MAPPO están en la Tabla~\ref{tab:plan2-hiper}.

El programa fijo se vuelve a tunear en cada celda con un barrido del
verde de la avenida, el verde transversal y el desfase entre cruces
sobre las semillas 2001 a 2003, por retraso medio. La malla del
barrido se amplía hasta que el óptimo queda en su interior; el verde
mínimo de 10~s es la única frontera que se acepta como óptimo. El
programa ganador se instala en el escenario de la celda y de él se
genera el actuado, de modo que los dos controles comparten fases.



La evaluación final corre cada modelo con la política congelada
sobre las semillas 1001 a 1030, las mismas treinta para todos los
brazos, y se evalúan las tres semillas base, no solo la publicada: el
rango entre ellas mide la fragilidad del entrenamiento.
La Tabla~\ref{tab:plan2-semillas} resume el uso de cada semilla.

\begin{table}[H]
\setstretch{1}
\centering
\footnotesize
\begin{tabular}{l>{\raggedright\arraybackslash}p{4.2cm}>{\raggedright\arraybackslash}p{6.2cm}}
\toprule
Etapa & Semillas de SUMO & Uso \\
\midrule
Tuneo & 2001 a 2003 & calibrar la demanda de P4 y barrer el programa fijo de cada celda \\
Entrenamiento & base$+k$, $k = 1 \ldots 200$, con bases 42, 2042 y 4042 (43 a 242, 2043 a 2242,
  4043 a 4242) & episodios de entrenamiento, con exploración \\
Validación & 999 y 1999 & cada 10 episodios y en el último: elige el \textit{checkpoint} publicado
  y la base publicada \\
Evaluación & 1001 a 1030 & política congelada, todos los brazos y todas las bases; única fuente
  de los veredictos \\
\bottomrule
\end{tabular}
\caption{Semillas de cada etapa del plan de pruebas. Los cuatro rangos no se solapan.}
\label{tab:plan2-semillas}
\end{table}


\subsection*{Métricas de salida}

Toda corrida, de entrenamiento o de evaluación y de cualquier modelo,
devuelve las mismas columnas. De los registros por vehículo de SUMO
salen el retraso (\texttt{timeLoss}), la espera detenida, las paradas
por vehículo, el \textit{throughput}, los vehículos no terminados y las
emisiones de CO$_2$ y NO$_x$; de TraCI, la cola media y máxima, las
decisiones, los cambios de fase y las métricas de corredor y
\textit{spillback} de la Sección~\ref{sec:u3}. Además se miden cada
segundo, durante todo el episodio y sobre todos los semáforos, tres
retornos, con suscripciones de TraCI para no multiplicar las
llamadas:
\begin{itemize}
  \item \texttt{ret\_retraso}: la suma, sobre los segundos del episodio
    y los semáforos, del término por segundo de la
    Ecuación~\ref{eq:recompensa-u4} con $w_3 = 0$,
    $-\bigl(w_1 \sum_a n_a (1 - \bar v_a / v_{\max,a}) + w_2 \sum_a
    \text{espera}_a\bigr)$. Es la recompensa de entrenamiento común,
    sin dividir entre 5 ni escalar.
  \item \texttt{ret\_colas}: la misma suma con la recompensa de U2
    (Ecuación~\ref{eq:recompensa}) evaluada cada segundo,
    $-\bigl(w_1 \sum_a c_a + w_2 \sum_a \text{espera}_a\bigr)$.
  \item \texttt{ret\_spill}: \texttt{ret\_retraso} menos
    $w_3 \sum_t \sum_e \max\bigl(0,\, q_e(t) - \lfloor P_e/2 \rfloor\bigr)$
    sobre todos los enlaces entre semáforos, con $w_3 = 2$.
\end{itemize}
Se calculan igual para el fijo, el actuado y los aprendices, de modo
que cada modelo puede compararse con los demás en su propia recompensa
y no solo en el retraso. En la tabla de veredicto, la columna
$\Delta$~recompensa compara la media del retorno propio de cada modelo
(\texttt{ret\_retraso} para Q-learning, SARSA, DQN, IPPO \emph{local},
IPPO \emph{vecinos} y MAPPO; \texttt{ret\_colas} para
\texttt{ql\_colas}; \texttt{ret\_spill} para IPPO \emph{spill}) con la
del mismo retorno medido bajo el fijo, sobre las mismas semillas. Como
los retornos son negativos, el cambio se divide entre el valor absoluto
de la media del fijo: un valor positivo indica que el modelo obtiene
más recompensa que el fijo. Los retornos no entran en el veredicto.
% fuente: rl_ippo.py y rl_corredor.py (correr_episodio: ret_retraso, ret_colas, ret_spill);
% consolidar.py (delta sobre |media_ref|); veredicto.py (columna Delta recompensa)

Por brazo y métrica se reportan la media, la desviación estándar y el
intervalo de confianza al 95\,\% sobre las treinta semillas. Frente al
fijo y frente al actuado se reportan el cambio porcentual de la media,
su intervalo al 95\,\% por \textit{bootstrap} pareado por semilla
(10\,000 remuestreos) y el valor $p$ de la prueba de Wilcoxon pareada.
El emparejamiento es por semilla, no por posición en el archivo. En
los retornos, que son negativos, el cambio se calcula sobre el valor
absoluto de la referencia, de modo que un cambio positivo es una mejor
recompensa. La métrica primaria es el retraso por vehículo.

Las llegadas son equiespaciadas en todos los casos
(Sección~\ref{sec:escenarios}): la variabilidad entre semillas
proviene solo del comportamiento de los conductores, por lo que los
intervalos de confianza son estrechos. Las llegadas aleatorias,
con un proceso de Poisson, cambian también los controles de referencia
y quedan fuera de este plan. Con la demanda nominal todos los
vehículos terminan su viaje bajo los controles de referencia y solo
los que entran al final de la hora llegan después de los 3600~s, así
que el \textit{throughput} y los vehículos no terminados discriminan
solo en P2. En P4, con $f = 1.2$, el actuado termina todos los
vehículos en las semillas de calibración
(Tabla~\ref{tab:red-alta-factor}).


\subsection*{Regla de veredicto}

El veredicto de un modelo en una celda se decide con el cambio del
retraso de su política publicada frente al fijo ($\Delta_f$, con su
intervalo y su valor $p$) y frente al actuado ($\Delta_a$):
\begin{itemize}
  \item \textbf{FUNCIONA:} el intervalo de $\Delta_f$ queda entero por
    debajo de cero, $p < 0.05$, $\Delta_f \leq -5$\,\% y el modelo no
    pierde contra el actuado de forma significativa (intervalo de
    $\Delta_a$ entero por encima de cero con $p < 0.05$).
  \item \textbf{MEJORA:} el intervalo de $\Delta_f$ queda entero por
    debajo de cero y $p < 0.05$, pero la mejora es menor del 5\,\% o el
    modelo pierde contra el actuado de forma significativa.
  \item \textbf{FALLA:} el intervalo de $\Delta_f$ queda entero por
    encima de cero y $p < 0.05$.
  \item \textbf{EMPATA:} cualquier otro caso.
  \item \textbf{PENDIENTE:} la celda no está completa o la base
    publicada no convergió; en el segundo caso se escribe también el
    veredicto provisional.
\end{itemize}
El umbral del 5\,\% separa una diferencia estadísticamente clara pero
pequeña de una que justifica cambiar de controlador; la meta de la
tesis, $-25$\,\%, es más exigente y se contrasta aparte. La
convergencia se mide sobre la serie de validación de la base
publicada, en el último tercio de las validaciones. La meseta es la
pendiente de la regresión lineal multiplicada por un tercio de los
episodios y dividida entre la media del tercio, en porcentaje: cuánto
cambiaría el valor si la tendencia siguiera un tercio más de
entrenamiento. La estabilidad es la distancia porcentual entre la
media de las tres últimas validaciones y el mínimo de la media móvil
de tres. La política converge si la meseta está entre $-5$ y
$+5$\,\% y la estabilidad no pasa del 5\,\%. Se reporta también
cuántas de las tres bases convergen.

El veredicto sale solo del retraso. La recompensa de cada modelo es
un diagnóstico: dice si el aprendiz optimizó lo que se le pidió, no si
controla mejor el tráfico, y dos recompensas razonables pueden ordenar
a los controles al revés. En corredor\_alta, con las diez semillas de
evaluación de la fase anterior, el actuado baja el retraso un 20.1\,\%
frente al fijo (intervalo de $-22.0$ a $-18.0$\,\%) y a la vez empeora
un 28.1\,\% la recompensa de colas de U2 por decisión ($-30.60$ frente
a $-23.89$); por la recompensa de colas se diría que el actuado es
peor, y por la métrica que la tesis se propone reducir, que es mucho
mejor.
% fuente: evaluacion_corredor_alta_resumen.csv (actuado: timeloss_prom y recompensa_decision)

Cada modelo tiene una tabla de veredicto con una fila por caso y estas
columnas: \emph{Caso} y \emph{Escenario}; \emph{Sem.}, el número de
semáforos (con la rotonda cuando la hay); \emph{Carriles}, por sentido
en la avenida; \emph{Demanda}, en veh/h durante la hora simulada;
\emph{Retraso}, la media por vehículo de la política publicada con su
intervalo al 95\,\% sobre las semillas 1001 a 1030; \emph{$\Delta$
retraso vs fijo} y \emph{$\Delta$ retraso vs actuado}, el cambio
porcentual con su intervalo \textit{bootstrap} y el $p$ de Wilcoxon;
\emph{$\Delta$ recompensa}, el cambio de la recompensa propia del
modelo (retraso, colas o \textit{spillback}, según la
Tabla~\ref{tab:plan2-modelos}) frente a esa misma recompensa medida en
el fijo, positivo si mejora; \emph{Convergió}, si la política publicada
cumple el criterio de convergencia y, entre paréntesis, cuántas de las
tres bases lo cumplen; y \emph{Veredicto}, la categoría de la regla. El
actuado aparece como fila de referencia, juzgado frente al fijo. Una
matriz de caso por modelo resume todos los veredictos, y una tabla por
modelo lista las tres bases de cada caso para medir la fragilidad.

El script \texttt{veredicto.py} genera, a partir de las evaluaciones
consolidadas, las tablas de veredicto, la matriz y las tablas de
fragilidad, que se reportan en el Capítulo~\ref{cap:resultados}.
% Las tablas estan en secciones/tablas/veredicto_*.tex y se incluyen desde capitulo5.tex.
